%%%
%
% $Autor: Wings $
% $Datum: 2021-05-14 $
% $Pfad: GitLab/MLEdgeComputer $
% $Dateiname: AdapterBoard
% $Version: 4620 $
%
% !TeX spellcheck = de_DE/GB
% !TeX program = pdflatex
% !BIB program = biber/bibtex
% !TeX encoding = utf8
%
%%%



\chapter{Configuration of the Arduino IDE for the Arduino PortentaH7}

\section{Introduction}

In this chapter, the process of connecting the Arduino Portenta H7 to a computer or laptop and configuring essential settings to get started is explored. The steps for installing the necessary tools and ensuring the board is correctly set up in the Arduino IDE are detailed. Finally, a simple example sketch is demonstrated to verify that the board is working as expected and ready for further development.


There are set of examples which are build in Arduino (IDE) for the testing purpose, for checking all the configuration and setting up the board we can open one of the basic example \FILE{blink.ino} first.

\subsection{Configuration for the Arduino Portenta H7}

To program the \textbf{Arduino Portenta H7} in an offline environment, follow these steps:


\begin{enumerate}
	\subsection{Installation of the driver}
	\item \textbf{Installation of the driver:} Begin by installing the latest version of the Arduino IDE on your computer. Refer to Section \ref{sec:Installation of the Arduino IDE} for more details.
	\item \textbf{Installation of the packages:} Once the IDE installation is complete, open the \menu{Arduino IDE} and navigate to the menu  \menu{Tools}, located in the upper-left corner.
	\item In the menu \menu{Tools}, select the \menu{Board Manager}.
	\item In the window \menu{Board Manager}, use the search bar to locate the \textbf{Arduino Portenta H7} by typing its name as shown in Figure \ref{fig:Arduinoportentaboard}		.
	\item From the search results, select \menu{Arduino Mbed OS Portanta Boards} and click \menu{Install} to install the required package.
\end{enumerate}

\begin{center}
	\begin{tikzpicture}
		\node (G) at (0,0) {\includegraphics[width=8cm]{Arduino/ArdiunoIDE2/portentaboard.png}};
		
		\draw[red,domain=0:360,line width=1.4pt] plot ({-3+1*cos(\x)},{0.7
			+0.5*sin(\x)});
		
		
		
		
	\end{tikzpicture}
	\captionof{figure}{Arduino Mbed OS Portenta Boards Installation}\label{fig:Arduinoportentaboard}
\end{center}	


\subsection{Connection and Configuration of the Arduino Board to the Computer}

To run either a built-in example or a custom program on the Arduino Portenta H7, follow these initial steps:

\begin{enumerate}
	
	\item Connect the Arduino Portentaboard to the computer using a USB-A-USB-micro-cable.
	\item Open the Arduino IDE on the computer. A blank environment page will appear, showing the default \PYTHON{void setup()} and \PYTHON
	{void loop()} functions.
	\item Navigate to the menu \menu{Tools}, select \menu{Board}, and choose the connected board, which is the \textbf{Arduino Portenta H7}, as shown in Figure~\ref{}.
	\Mynote{put here a figure[screenshot of the connected board Arduino Portenta H7]}
	
\end{enumerate}

This setup prepares the Arduino IDE to recognize and communicate with the board Arduino Portenta H7.



\subsection{Select the Appropriate Port}


After selecting the Arduino Portenta H7 board, the next step is to verify the connected port. To do this, the Arduino board must be set into Bootloader mode by pressing the white reset button on the board, as shown in Figure ~\ref{fig:ArduinoPortentaH7resetButton}.


\begin{center}
	\begin{tikzpicture}
		\node (G) at (0,0) {\includegraphics[width=8cm]{Arduino/ArdiunoIDE2/resetbutton.png}};
		
		\draw[green,domain=0:360,line width=1.4pt] plot ({-1.5+0.4*cos(\x)},{-1.6
			+0.4*sin(\x)});
	\end{tikzpicture}
	\captionof{figure}{Arduino Portenta H7 Reset Button}\label{fig:ArduinoPortentaH7resetButton}		
\end{center}

By pressing the white reset button, the Arduino board will enter Bootloader mode. It is important to verify that the orange Builtin-LED is illuminated, as shown in Figure ~\ref{fig:ArduinoPortentaH7ResetButton}.		

\Mynote{board Portenta H7}



\begin{center}
	\begin{tikzpicture}
		\node (G) at (0,0) {\includegraphics[width=8cm]{Arduino/ArdiunoIDE2/resetbutton.png}};
		
		\draw[green,domain=0:360,line width=1.4pt] plot ({-2.8+0.4*cos(\x)},{-0.8
			+0.4*sin(\x)});
		
		\draw[orange,domain=0:360,line width=1.4pt] plot ({-2.8+0.4*cos(\x)},{-2.6
			+0.4*sin(\x)});
		
		
		
	\end{tikzpicture}
	\captionof{figure}{green and orange Builtin-LED}\label{fig:ArduinoPortentaH7ResetButton}	
\end{center}	

\Mynote{board Portenta H7}



After successfully completing the previously mentioned steps, the next task is to select the connected port before uploading the program. To do this, navigate to the menu \menu{Tools}, select \menu{Port}, and ensure that the available port for uploading the program is properly selected, as shown in Figure ...


\Mynote{here put new figure [screenshot] available Port for uploading the sketch  - follow the steps above and change the screenshots with Arduino Portenta H7}	


\begin{figure}[H]\centering
	\includegraphics[width=8cm]{Arduino/ArdiunoIDE2/SelectAvailablePortForUploadingArduinoSketch}
	\caption{Select Available Port for Uploading Arduino Sketch}
	\label{fig:SelectAvailablePortforUploadingArduinoPortentaSketch}		
\end{figure}





\subsection{Upload and Verify a Sketch}\label{uploadcode}

After ensuring that the appropriate port is selected, the next step is to upload the Arduino program. \textbf{Before} uploading the program, it is considered best practice to verify it first. This process will indicate whether any errors or warnings exist in the program. Once the program is successfully verified, it can be safely uploaded by clicking the button \menu{Upload} located below the file section, as shown in Figure ~\ref{fig:Upload the Program in Arduino Portenta H7 board}




After selecting the board and port, confirm that the Arduino Portenta H7 board is properly connected to the Arduino IDE. This can be verified by checking the message in the bottom right corner of the IDE window, which should display the connected board and port.





\begin{figure}[H]\centering
	\includegraphics[width=8cm]{Arduino/ArdiunoIDE2/UploadTheProgramInArduinoBoard}
	\caption{Upload the Program in Arduino board}
	\label{fig:Upload the Program in Arduino Portenta H7 board}		
\end{figure}




After uploading, the code will be compiled, and any issues in the program will be displayed in the bottom black window. Once the code is successfully uploaded and compiled on the Arduino Portenta H7 board, it is necessary to reselect the port, as done previously. Navigate to the \menu{Tools menu}, select \menu{Port}, and ensure the correct port is selected, as shown in Figure~\ref{fig:SettingthePortH7}, in order to view the output in the Serial Monitor.



\begin{figure}[H]\centering
	\includegraphics[width=8cm]{Arduino/ArdiunoIDE2/SettingThePort}
	\caption{Setting the Port}
	\label{fig:SettingthePortH7}		
\end{figure} 


\Mynote{check the figure.. try the steps above and make a new screenshot and correct the figure 3.5 and 3.6}




\section{Test the Configuration}


In this subsection, the board's connection to the Arduino IDE is tested by running an Example Blink Sketch. This simple test lights up an onboard LED, verifying both the connection and the basic functionality of the Arduino Portenta H7 board. By the end of this section, confirmation is provided that the hardware and software are working together seamlessly.

\subsection{Testing the Steps by an Example Sketch \FILE{blink.ino}}


The Arduino IDE contains a set of built-in examples for testing purposes. To verify the configuration and set up the board, the basic example \FILE{blink.ino} can be opened first. After uploading the example \FILE{blink.ino} the Builtin LED blinks with 2 Hz. 


To begin, open the Arduino IDE on the computer and follow the steps below:

\begin{enumerate}
	
	
	
	\item \textbf{Open the Examples:} Once it's open, click on the \menu{File} menu and hover over \menu{Examples} in the dropdown menu. 
	
	\item \textbf{Selecting the Example Sketch:} A list of example categories will appear. Under the Built-in Examples, go to \menu{01.Basics} and click on \menu{Blink} as shown in ~\ref{fig:LED-ExampleTestH7}.
	
	\item \textbf{Blink Example Sketch:} After following steps 1 and 2, the  Example Sketch \FILE{blink.ino} will open in a new window within the IDE as seen in Figure ~\ref{fig:LED-Built_CodeH7}.
	
	\item \textbf{Upload Example Sketch \FILE{blink.ino}:} Click the \menu{Upload} button (right arrow icon) in the Arduino IDE toolbar to compile and upload the example sketch to the Arduino Portenta H7 board. Once the upload is complete, the orange built-in LED starts blinking. 
	
	
	
\end{enumerate}


\begin{center}
	\includegraphics[width=8cm]{Arduino/ArdiunoIDE2/blink_example_IDE2.3.3.png}
	\captionof{figure}{LED Example Sketch}\label{fig:LED-ExampleTestH7}		
\end{center} 



\begin{center}
	\includegraphics[width=8cm]{Arduino/ArdiunoIDE2/ledblink_code.png}
	\captionof{figure}{LED-Built Example}\label{fig:LED-Built_CodeH7}		
\end{center} 

With the successful blinking of the built-in orange LED, it is confirmed that the Arduino Portenta H7 is properly connected to the Arduino IDE and its basic functionality has been verified. The board is now ready for further development and projects.


\subsection{Example ``Hello World'' - Sketch Blink}

To upload the example sketch, go to \SHELL{File > Examples > Basics > Blink} as shown in Figure ~\ref{UploadSketch}. \\

\begin{figure}
	\begin{center}
		\includegraphics[width=0.7\linewidth]{Arduino/PortentaH7/UploadSketch.png}
		\caption{Upload Basic Sketch}
		\label{UploadSketch}
	\end{center}
\end{figure}

In the program, we initialize \SHELL{LED\_BUILTIN} as the output in the setup function. In the loop function we turn on the LED by using \SHELL{digitalWrite()} function. Click on upload and wait for it to complete. After the upload is done, the green LED on the board starts blinking with a delay of 1000ms. The blink sketch appears on the screen as shown in figure \ref{CompileSketch}.


{
	
  \captionof{figure}{Program ``Hello World'' - \FILE{Blink.ino}}
  
  \ArduinoExternalO{../../Code/ArduinoCLI/Blink.ino}
}




	\begin{center}
		\includegraphics[width=0.7\linewidth]{Arduino/PortentaH7/CompileSketch.png}
		\captionof{figure}{Compile Blink Sketch}
		\label{CompileSketch}
	\end{center}


\subsection{Serial Blink Example for Dual Core Processing}

The following figure \ref{SerialBlinkSketch} represents the serial blink sketch. This sketch helps in demonstrating the serial data transmission by showing the LED blinking after a certain time period. \\

Here, \SHELL{Serial.begin(115200)} command sets the data transmission rate to 115200 bits per second. The LEDB is the output which blinks in blue color and we have specified core M7. If we would like to use core M4 then we need to boot it first using bootM4() command. The ouput LED blinks blue in color after a delay of 3 seconds as defined in the sketch. \\


In order to use core M4, we need to de-select the core M7 from the board and select core M4. After selecting, we can see at the bottom the selected core for confirmation. Also, We need to write another sketch and now change the LED color to red in the sketch and add a delay of 4 seconds. 



	\begin{center}
		\includegraphics[width=0.7\linewidth]{Arduino/PortentaH7/SerialBlink.png}
		\captionof{figure}{Serial Blink Sketch}
		\label{SerialBlinkSketch}
	\end{center}



Notice that in the output both LED blue and green blinks and at a point in time they almost blink together resulting in a Lila color. This shows the dual core processing concept by using an LED. \\

\begin{figure}
	\begin{center}
		\includegraphics[width=0.7\linewidth]{Arduino/PortentaH7/Output.png}
		\caption{Output LED Color}
		\label{Output}
	\end{center}
\end{figure}

It can be seen from the figure \ref{Output} that the output LED blinks in Lila color due to the overlapping of the green and blue LED blinking at the same time.



	\subsection{Example Blink Sketch}

To upload the example sketch, go to File then click on examples and then click on 01.Basics and then select blink as shown in Figure \ref{fig:Example LED-Test}.

\begin{center}
	\label{fig:Arduino PortentaH7 Connected to a Laptop}
	\includegraphics[width=0.7\linewidth]{Arduino/PortentaH7/Connectedtoalaptop.png}
	\captionof{figure}{Arduino PortentaH7 Connected to a Laptop}
\end{center}

\begin{center}
	\label{fig:Example LED-Test}
	\includegraphics[width=0.7\linewidth]{Arduino/PortentaH7/MenuBarOptions.png}
	\captionof{figure}{Example LED-Test}
\end{center}


In the program, we initialize LED\_BUILTIN as the output in the setup function. In the loop function we turn on the LED by using digitalWrite() function. Click on upload and wait for it to complete. After the upload is done, the green LED on the board starts blinking with a delay of 1000ms.The blink sketch appears on the screen as shown in Figure \ref{fig:Blink Sketch Compile}.

\begin{center}
	\label{fig:Blink Sketch Compile}
	\includegraphics[width=0.7\linewidth]{Arduino/PortentaH7/CompileBlinkSketch.png}
	\captionof{figure}{Blink Sketch Compile}
\end{center}


\subsection{To-Do's}


\begin{enumerate}

\item  compare Configuration of the Arduino Portenta H7 and Configuration of the Arduino BLE Sense 

\item check the Mynote comemnts and refresh the figures.. follow the steps with an Arduino Portenta H7 board and make new screenshots on the right posositions (check the comments on the side) 

\item


\end{enumerate}

